\begin{center}\textbf{27.}\end{center}
\noindent\hspace*{\parindent}Neben diesen nur handschriftlich vorhandenen Sätzen stehen die in den genannten Artikeln der D. A. angeführten asymptotischen Bestimmungen des dort definierten »Mittelwerts« für die Anzahl der Geschlechter eigentlich primitiver binärer quadratischer Formen als Funktion ihrer Determinante \(\pm D\):
\[
\frac4{\pi^2}\mathbin{.}\log D+\frac8{\pi^2}C+\frac{48}{\pi^4}\mathbin{.}\sum\frac{\log n}{n^2}-\frac{2\log2}{3\pi^2},
\]
wo \(C\) die \textsc{Euler}–\textsc{Mascheroni}sche Konstante bedeutet, und des Mittelwerts für die Anzahl ihrer Klassen bei negativer Determinante \(-D\):
\[
\frac{2\pi}{7\mathbin{.}\displaystyle\sum\frac1{n^3}}\mathbin{.}\sqrt D-\frac2{\pi^2}.
\]
Obwohl auch hier jede Andeutung der angewandten Methode fehlt, wird doch ausdrücklich von \textsc{Gauss} ausgesprochen (art. 304), dass er durch eine theoretische Untersuchung, deren Grundsätze er bei anderer Gelegenheit mitteilen wolle, zu diesen Formeln gelangt sei, und ein schriftlicher Vermerk (W. I, S. 476, Anmerkung zum art. 302) in seinem Handexemplare der D. A. setzt die Entstehung derselben auf den Anfang des Jahres 1799. In der Tat findet man in \textsc{Gauss}’ im November 1799 begonnenem Heftchen (Ac) S. 19 (abgedruckt W. X1, S. 15) die erstere Anzahl, freilich erst in der unbestimmteren Form
\newpage\SourceLeaf{74}\ArithmeticHead{67}{R}{MITTLERE WERTE. BESTIMMUNG DER KLASSENANZAHL.}\HistoricalFoot{}{9*}
\[
a\log D+\beta,
\]
und an anderer Stelle (Eb) schon etwas bestimmter und mit dem Zusatz: »Multitudo generum per disquisitionem theoreticam inventa est« durch
\[
\frac8{\pi^2}\log D+\mathrm{const.}\footnotemark
\]
\footnotetext{Im Nachlasse steht in diesen Formeln \(N\) statt \(D\), sicher aber in gleicher Bedeutung.}
ausgesprochen. Der für die zweite Anzahl angegebene Ausdruck, abgesehen von der Konstanten \(-\frac2{\pi^2}\), welche \textsc{Gauss} auch an späterer Stelle (W. II, S. 284) unterdrückt hat, tritt bereits in \textsc{Gauss}’ im November 1798 begonnenem Heft (Ab) S. 8 auf, ohne Angabe seiner Bedeutung. Die entsprechende Untersuchung für Formen mit positiver Determinante, deren Resultat schon in art. 304 D. A. richtig vermutet wird, erledigte sich nach dem Additamentum ad. art. 306 X (W. I, S. 466) und der Anmerkung zu demselben (ebendas. S. 476) erst am Ende des Jahres 1800. In dieser Zwischenzeit muss \textsc{Gauss} die analytischen Betrachtungen ausgebildet haben, die er in zwei Abhandlungen aus den Jahren 1834, 1837 niedergelegt hat und die (W. II, S. 269) mit ausführlichen Erläuterungen \textsc{Dedekinds} abgedruckt sind, Betrachtungen, welche ihn schon vor \textsc{Dirichlet} zur Bestimmung des wahren Wertes der Klassenanzahl geführt haben.

Sie gründen sich auf denselben geometrischen Hilfssatz wie bei \textsc{Dirichlet}. Denkt man sich im rechtwinkligen Gitter ganzzahliger Punkte, wie es bei den komplexen ganzen Zahlen \(a+bi\) auftritt, eine geschlossene Kurve vom Inhalte \(V\) und nennt \(M\) die Anzahl der Gitterpunkte (oder Quadrate \(\kappa^2\)), welche in ihr Inneres fallen, so ist, wenn die Kurve gleichmässig nach allen Richtungen unendlich erweitert wird, \(\lim\frac MV=1\), oder anders ausgedrückt: bleibt \(V\) konstant, nehmen aber die Quadrate unendlich ab, so ist
\[
\lim_{\kappa=0}\frac{M\kappa^2}{V}=1.
\]
Da die Anzahl der Punkte innerhalb der Kurve
\[
ax^2+2bxy+cy^2=A
\]
bei negativer Determinante \(-D\) die Anzahl \(F(A)\) der Darstellungen aller Zahlen \(1,2,3,\ldots A\) mittels der Form \((a,b,c)\) ist, so ergibt sich dem Hilfs-
\newpage\SourceLeaf{75}\ArithmeticHead{68}{L}{BACHMANN, ÜBER GAUSS’ ZAHLENTHEORETISCHE ARBEITEN.}
\noindent satze zufolge für die mittlere Anzahl der Darstellungen einer Zahl \(A\) durch dieselbe Form die Gleichung
\[
\lim_{A=\infty}\frac{F(A)}A=\frac\pi{\sqrt D}.
\]
Da sie nur von der Determinante \(D\) abhängt, so erschliesst man für die mittlere Anzahl der Darstellungen einer Zahl \(A\) durch das gesamte Formensystem mit der Determinante \(-D\) den Wert
\[
h\mathbin{.}\frac\pi{\sqrt D},
\]
in welchem \(h\) die Klassenzahl dieses Systems bedeutet.

Stimmen soweit die Methoden beider Forscher mit einander überein, so unterscheiden sie sich in der Herleitung eines zweiten Grenzwertes für dieselbe mittlere Anzahl d. i. in deren Zurückführung auf die Summe
\[
\sum_{n=1}^{\infty}\left(\frac Dn\right)\frac1n,
\]
bezüglich welcher auch die erhaltenen \textsc{Gauss}schen Notizen nicht genügenden Aufschluss geben. An einem besonders glücklichen Tage aber, dem 30. November 1800 (Nr. 114 T.), gelang es \textsc{Gauss}, die Klassenanzahl für Formen mit negativer Determinante unter dreifacher Gestalt zu bestimmen: einmal mittels der gedachten Summe, ein zweites Mal unter der Gestalt eines unendlichen Produktes, ein drittes Mal mittels trigonometrischer Funktionen (s. W. II, S. 285, wo auch für die Klassenanzahl bei positiver Determinante entsprechende Ausdrücke gegeben sind), und er fand endlich noch am 3. Dezember desselben Jahres (Nr. 115 T.) für Formen mit negativer Determinante \(-D\) jenen einfachen, nur aus der Menge der Zahlen unterhalb \(D\), welche Teiler bezw. Nichtteiler der Form \(x^2+D\) sind, gebildeten Ausdruck (s. ihn W. II, S. 286, [VII] für den Fall, wo \(D\) eine Primzahl von der Form \(4n+1\) ist), den für diesen einfachsten Fall einer Primzahldeterminante \(-p\) zuerst \textsc{Jacobi}\footnote{\textsc{Crelles} Journal für Mathematik \textbf{9} (1832), S. 189, Werke VI, S. 420.} bekannt gemacht hat.

Mit Hilfe dieses Ausdrucks konnte \textsc{Gauss} eine Reihe von Sätzen über die
\newpage\SourceLeaf{76}\ArithmeticHead{69}{R}{BESTIMMUNG DER KLASSENANZAHL.}
\noindent Verteilung der quadratischen Reste \((\mathrm{mod.}\ p)\) auf die Achtel bezw. Zwölftel des Intervalls \(1\) bis \(p\) aufstellen (W. II, S. 288—291, 301—303), sowie endlich auch, was gelegentlich der \textit{Comm. prima} über biquadratische Reste noch als eine Lücke angedeutet worden, für den Fall \(p\equiv5\ (\mathrm{mod.}\ 8)\) die, für den andern Fall \(p\equiv1\ (\mathrm{mod.}\ 8)\) bis auf den heutigen Tag noch ungelöste Frage entscheiden (W. II, S. 287), mit welchem Vorzeichen behaftet die Basis \(b\) der Zerlegung \(p=a^2+b^2\) im Schlusstheoreme jener Abhandlung erhalten werde. Diese Verknüpfung einander scheinbar sehr entlegener Gegenstände und Gedanken, hat \textsc{Gauss} selbst als eine sehr verschlungene und interessante bezeichnet (Brief an \textsc{Dirichlet} vom 30. Mai 1828, W. II, S. 516).

\bigskip
So sehen wir, wohin wir uns auch wenden, \textsc{Gauss}’ Forschung schon in jedes der verschiedenen Gebiete, welche die heutige Wissenschaft der höheren Arithmetik umfasst, tief und bahnbrechend eindringen; keins, das nicht — die additive Zahlentheorie etwa ausgenommen — auf \textsc{Gauss}’ Arbeiten als seine Grundlage oder den Keim seiner Entwicklung zu weisen hätte. Denn auch die Idealtheorie der algebraischen Zahlenkörper weist uns auf die von \textsc{Gauss} erkannte Notwendigkeit der Zahlen \(a+bi\) zurück und ist nur ihr Analogon auf höherer Stufe.

\begin{center}\small SCHLUSSBEMERKUNG.\end{center}
\begingroup\small
\noindent\hspace*{\parindent}Die meisten textlichen Änderungen und neu hinzugekommenen Fussnoten, die der vorstehende Abdruck meines Aufsatzes gegenüber der ursprünglichen, 1911 erschienenen Fassung aufweist, sind aus Mitteilungen der Herren \textsc{F. Klein}, \textsc{L. Schlesinger} und \textsc{P. Stäckel} hervorgegangen. Herrn \textsc{Schlesinger} ist auch die übersichtliche Gliederung zu verdanken, die der Aufsatz erhalten hat.
\begin{flushright}\textsc{Bachmann.}\end{flushright}
\endgroup
\begin{center}\rule{.24\textwidth}{.3pt}\end{center}
